\section{Limitations}
\label{sec:limitations}

\textbf{Environments, devices, and seeds.} The study covers three benchmark environments and one external auditorium, recorded with two smartphone models. It is not a cross-device study, and RTT and RSS are known to depend on chipset, antenna, and firmware. The main tables use one seed; the headline comparison and the EETAC protocols were repeated with ten seeds, but the Building literature-baseline and selection-rule analyses use one seed because of their cost. At full radio-map density, the seed standard deviation of the paper-compatible metric is at most $0.008$~m, well below the differences between $\rho=0.7$ and $\rho=1$ but comparable to some differences among the randomized ensembles of Table~\ref{tab:literature_baselines}. The bootstrap intervals resample test RPs for a fixed model and do not include seed variability.

\textbf{The subspace ratio.} $\rho=0.7$ is a robust default, not the optimum: at full density, the best ratios are $0.5$--$0.85$, and on sparser maps they are smaller. The decrease of the best ratio with density is observed in all three environments, but for Random Forests only, and in Office (28 training RPs) the gain over bagging does not grow with sparsity; extremely randomized trees hardly benefit from feature subsampling at all. A rule that sets $\rho$ from radio-map density, AP count, or modality reliability is suggested by the results but not developed or validated.

\textbf{Development exposure.} The original benchmark was used during method development, so its official-test improvements are not results on a completely untouched benchmark. This is one reason for the external EETAC evaluation, for which $\rho$ was frozen beforehand. In the extension analyses, the official tests score alternatives in hindsight but are never used to choose the proposed configuration. Moreover, the Apartment improvement on the official test is not significant on its own (its bootstrap interval includes zero), although it occurs in all ten seeds.

\textbf{EETAC.} The mirrored frame of the second acquisition day is inferred from the data (22 of 25 RPs match the mirrored positions, and the identity matches only the mirror-invariant central row). The inference is strong but has not been confirmed by the dataset authors, and which day uses the original frame cannot be determined from the files; accuracy does not depend on this choice. The $5\times5$ grid is a sparse radio map, so RP-disjoint EETAC errors (2.4--2.9~m) are not representative of densely surveyed deployments, and the two-AP configuration of~\cite{gonzalezdiaz2026scalable} could not be reproduced unambiguously.

\textbf{Temporal robustness.} Temporal shift is evaluated with a single pair of acquisition days in one environment, and in one of the four cross-day comparisons the two models tie. Longer time spans, changes in occupancy or furniture, and AP replacement are not covered.

\textbf{Tail errors.} The configuration improves average metrics but not every tail statistic: the Apartment official-test $P_{90}$ and the EETAC sample-level $P_{90}$ are worse than for the full-feature forest. Applications with strict worst-case requirements need mechanisms that target the tail.

\textbf{Uncertainty.} The covariance-aware ellipse reduces area but does not improve coverage. The group calibration is valid only if RPs are exchangeable, which a regular survey grid satisfies only approximately; it discards all but one scan per RP and gives no finite region with fewer than $1/\alpha-1$ calibration RPs. The RP-count correction is a conservative heuristic without a finite-sample guarantee, and it degenerates to the largest calibration score when fewer than $2/\alpha-1$ RPs are available. The full-training grouped-OOF variant has no finite-sample guarantee. Locally normalized scores based on the spread of the tree predictions~\cite{lei2018distribution} reduced the area by $25\%$ at equal coverage in Apartment but over-covered in Building in a screening experiment and are left for future work.

\textbf{Efficiency and comparison.} Timing and model size are measured in Python on a laptop CPU; energy, peak memory, and latency on a phone or embedded device are not measured. Deep learning models are not re-implemented, and the comparison with GTCN relies on its published Building result~\cite{rana2025gtcn}, which is more accurate than the forests studied here. Claims are therefore limited to the experimentally verified comparisons.

\textbf{Feature importance.} The modality analysis relies on impurity-based RF importance, which describes the fitted models and can be redistributed among correlated features; it does not measure the intrinsic value of RTT relative to RSS.
